\section{Background}

Skin cancer is among the most common cancers worldwide, and the prognosis for a
patient depends heavily on how early the lesion is found. Melanoma in
particular is highly treatable when caught early but becomes far more dangerous
once it spreads. Dermoscopy --- a non-invasive imaging technique that removes
surface glare and reveals the sub-surface structure of a lesion --- has become a
standard tool for this kind of early assessment. Reading dermoscopic images
well, however, takes years of training, and even experienced dermatologists do
not always agree with one another. This gap between the volume of lesions that
need screening and the number of specialists available to read them is what
motivates automated, image-based diagnostic support.

Deep learning has reshaped this area over the past decade. Convolutional neural
networks (CNNs) trained on large dermoscopic collections reached
dermatologist-level performance on melanoma recognition~\cite{esteva2017dermatologist},
and the annual ISIC challenges have since turned skin lesion classification into
a well-defined benchmark task. The ISIC-2019 dataset, which we use throughout
this work, contains 25{,}331 labelled dermoscopic images spanning eight
diagnostic categories: melanoma (MEL), melanocytic nevus (NV), basal cell
carcinoma (BCC), actinic keratosis (AK), benign keratosis (BKL),
dermatofibroma (DF), vascular lesion (VASC), and squamous cell carcinoma (SCC).

More recently, Vision Transformers (ViTs)~\cite{dosovitskiy2021vit} and hybrid
designs that place a transformer on top of convolutional
features~\cite{mehta2022mobilevit,li2022efficientformer} have shown that
modelling the global structure of a lesion --- its overall asymmetry, the
regularity of its border --- complements the local texture that convolutions
capture so well. These two views of a lesion, the local and the global, are
exactly the cues a clinician is trained to weigh, which makes the
CNN--Transformer combination a natural fit for dermoscopy.

\section{Rationale of the Study or Motivation}

Two practical problems keep most published skin-lesion models from being used in
an everyday clinic. The first is cost. The strongest results on ISIC-2019 come
from large ensembles of heavy networks~\cite{isic2019leaderboard,gessert2020ensembles};
these are accurate but impractical to run on the modest hardware found in a
typical dermatology practice or on a portable device next to the patient. The
second is that image-only models ignore information a dermatologist always has
in front of them. The age of the patient, their sex, and the part of the body
where the lesion sits all shift the prior probability of each diagnosis --- a
scaly patch on the face of an elderly patient is read very differently from the
same patch on a young person's back. This clinical context is recorded in the
ISIC-2019 metadata, yet it is routinely discarded by classifiers that look only
at pixels.

Our study is driven by the belief that these two problems should be solved
together rather than separately. We argue first that a carefully designed
lightweight hybrid can match or beat much larger image-only baselines on
balanced accuracy while staying small enough to deploy, and second that the
clinical metadata a dermatologist already has on hand can be folded in at
negligible cost as additional diagnostic context, handled so that missing
fields are never imputed. We
also recognise that not every deployment has the same constraints, so we deliver
two models from one design philosophy: a genuinely compact model for
constrained settings, and a higher-capacity model for situations where balanced
accuracy matters more than footprint.

\section{Problem Statement}

ISIC-2019 is severely imbalanced. The most common class, nevus, has roughly
12{,}875 images, while dermatofibroma and vascular lesions have only about 239
and 253 respectively --- a ratio of more than fifty to one. A model trained
naively on this data learns to predict the common classes and quietly ignores
the rare ones, which is the opposite of what matters clinically, since several
of the rare classes are malignant. Reporting overall accuracy hides this failure
entirely. The problem we address is therefore to build a model that performs
well \emph{across all eight classes at once} --- measured by balanced
multi-class accuracy (BMA), the mean of the per-class recalls --- while
respecting a tight parameter and compute budget and while making use of the
clinical metadata that accompanies each image, including the large fraction of
records where some of that metadata is missing.

\section{Objective}

The objectives of this work are:

\begin{enumerate}
  \item To design a lightweight hybrid CNN--Transformer classifier for
        ISIC-2019 that stays under 6M parameters and 1 GMAC, and that fuses
        clinical metadata through a cross-attention head capable of handling
        missing fields without imputation.
  \item To design a second, higher-capacity model in the same family --- a
        dual-backbone network with learned attention fusion and a
        Kolmogorov--Arnold Network (KAN) classifier --- that trades footprint
        for a higher balanced accuracy.
  \item To train both models, together with five established lightweight
        baselines, under one identical loss, augmentation, and evaluation
        pipeline so that the comparison is fair.
  \item To quantify the contribution of each design choice (the CNN stem, the
        transformer trunk, the metadata head, and the dual-backbone fusion)
        through controlled ablations, and to report efficiency alongside
        accuracy.
\end{enumerate}

\section{Methodology in Brief}

Every image passes through a fixed preprocessing pipeline: near-duplicate
removal by perceptual hashing, a lesion-grouped stratified train/validation/test
split, Shades-of-Gray colour constancy~\cite{finlayson2004shadesofgray}, and
resizing to $224\times224$. The cleaned images are stored as a single memory-mapped
array so that training does not pay a per-batch decoding cost. All models are
trained with Class-Balanced Focal Loss~\cite{cui2019classbalanced,lin2017focal}
to counter the imbalance, with AdamW, a cosine schedule with warmup, weight
EMA, and eight-view test-time augmentation at inference. The lightweight model
combines a truncated MobileNetV2 stem with a small ViT trunk and a metadata
cross-attention head; the flagship model (DEKAN) fuses DenseNet-121 and
EfficientNet-B0 features through learned attention before a transformer trunk and
a KAN classifier. The primary metric throughout is BMA.

\section{Scopes and Challenges}

The scope of this thesis is the eight-class ISIC-2019 classification task and
the two proposed architectures, evaluated against same-budget lightweight
baselines on a held-out test split. The main challenges were the extreme class
imbalance, the substantial missingness in the metadata (around 30\% of age
values, for example), and a hard hardware budget --- all experiments run on a
single 16~GB GPU under Windows, which rules out multi-worker data loading and
forces aggressive use of mixed precision and a memory-mapped dataset. We did not
extend the study to external test sets or to multi-seed reporting; both are
discussed openly as limitations in Chapter~5.

\section{Thesis Organization}

Chapter~2 reviews the architectures and prior work this study builds on.
Chapter~3 states the data, requirements, constraints, and broader impact of the
system. Chapter~4 describes the preprocessing pipeline, the two proposed models,
the loss, and the training and evaluation protocol. Chapter~5 presents and
analyses the results, including ablations, efficiency, and a comparison with
prior work. Chapter~6 concludes and outlines future work.
